% ── CHAPTER 9: CONVERGENCE 3: OBSERVER-DEPENDENCE (~6,000 words) ──
\chapter[Observer-Dependence]{Observer-Dependence: What You Can See Depends on What You Are}
\label{ch:observer-dependence}

% \epigraph{There is no view from nowhere.}{Thomas Nagel}

\firstword{T}{homas} Nagel's phrase has become something of a philosophical catchword, but it earns its keep.
There is a persistent temptation to imagine that reality, properly described, is described from no particular place --- that the truest account is the view an omniscient being would have, standing outside the universe, seeing everything at once.
In this imagined God's-eye view, every object has determinate properties; every event has a definite character; every fact is just a fact, independent of who is looking or what they are capable of registering.

The three traditions I am tracing in this book all reject that picture.
Not because they are anti-realist.
Not because they think the world is just whatever we say it is.
But because they have each, through their own methods and their own categories, arrived at the same structural conclusion: the way reality presents itself is inseparable from the kind of thing doing the looking.
The properties of a quantum system are irreducibly relational.
The depth of an actual occasion's experience is a function of its own relational complexity.
The kingdom of existence a creature inhabits determines what aspects of reality are available to it and what aspects remain entirely beyond its reach.

This is not relativism.
A thermometer's reading is genuinely real.
A person's reading is also genuinely real.
They are not the same reading.
Neither is false.
Neither is the complete truth.
Reality is one; the faces it presents to its various interlocutors depend on what those interlocutors are.

This chapter takes up that theme --- observer-dependence, hierarchy of comprehension, and what it means to say there is no view from nowhere --- across all three traditions.
The convergence here is subtler than the bare claim that observation affects what is observed.
It is the stronger claim that what becomes definite, what is even in principle available to a given type of observer, is a structural function of the observer's own relational constitution.
Not just that the measurement disturbs the system.
But that what is real for one kind of being is unavailable to another --- not because of ignorance, not because of imprecision, but because reality's architecture is hierarchical and perspectival all the way down.

\section{The Physics: No Absolute State, Only Relative States}
\label{sec:ch9-physics}

In 1996, the physicist Carlo Rovelli published a paper that has since become one of the most discussed and debated contributions to the foundations of quantum mechanics~\citep{Rovelli1996}.
The paper's title is ``Relational Quantum Mechanics,'' and its central claim is as plain as it is radical: quantum states are not absolute.
A quantum system does not have a definite quantum state, period.
It has a quantum state relative to some other system.
Change the relational context, and you change the state.

This sounds paradoxical until you remember that the same move is perfectly familiar in a different domain.
In special relativity, there is no such thing as an object's velocity, period.
There is only an object's velocity relative to some reference frame.
An electron at rest in my laboratory frame is moving at a substantial fraction of the speed of light in the frame of a muon created by a cosmic ray.
Neither description is wrong, and neither is the ``real'' one.
Velocity is frame-relative, not frame-independent.

Rovelli's proposal is that quantum states work the same way.
A particle's spin is not spin-up or spin-down.
It is spin-up relative to this detector, spin-down relative to that one.
A photon is not polarized along this axis or that one.
It is polarized along a given axis relative to a given polarizer.
There is no quantum state that the system simply has, independent of any relational context.

The radicalism of this move becomes clear when you push on it.
In ordinary quantum mechanics, the measurement problem asks: when does the wavefunction collapse?
Before measurement, the system is in a superposition of many possible states.
After measurement, it is in one definite state.
But when exactly does the collapse happen?

Relational quantum mechanics dissolves the puzzle by recasting it.
There is no measurement-independent state that collapses at some special moment.
There is only the state relative to a given observer.
When that observer interacts with the system, a definite outcome becomes actual for that observer.
But a second observer, who has not yet interacted with either the first observer or the system, may legitimately and correctly describe both the first observer and the system as still in superposition.

This sounds like a contradiction until you notice that the two observers are not describing the same thing.
The first describes the system relative to herself; the second describes the system-plus-first-observer relative to himself.
Only when the second observer acquires information from the first does a definite outcome become actual for him.

The key phrase is ``acquires information.''
In relational quantum mechanics, the event of one system acquiring information about another is itself a physical event.
It leaves a physical trace.
And that trace is what constitutes an outcome as definite for the acquiring system.
The measurement event does not disturb a pre-existing, determinate state.
It constitutes the outcome as definite, within that specific relational context.
The measurement is constitutive, not merely revelatory.

This is the feature I want to dwell on, because it will recur in the process-philosophical and \bahai accounts.
In a naive picture of measurement, the system has a definite state, and the measurement reads off what that state is.
In relational quantum mechanics, the system acquires a definite state only in the context of a particular relational interaction.
The outcome is not discovered; it is, in a precise sense, created by the relational event.
Created, and yet genuinely real: not invented, not arbitrary, but constituted.

A homely image helps here: the difference between reading a lock and fitting a key.
A pre-existing description would say: the key reveals the lock's shape.
A relational description says: the key and the lock co-determine a fitting or non-fitting through their interaction.
The fitting is real.
But it belongs to neither the key nor the lock in isolation.

Let me add one clarification that will matter for the convergence.
Relational quantum mechanics is not claiming that the world is subjective, or that different observers are trapped in incompatible and irreconcilable realities.
The constraint is consistency: if two observers have interacted (have entered into a relational context with each other), their descriptions of the shared outcomes must agree.
Rovelli is explicit that the relational facts are objective facts --- they are just facts about relations rather than facts about isolated substances.
The spin of an electron relative to a detector is a perfectly objective fact.
It is just not a fact about the electron alone; it is a fact about the electron-detector interaction.

This has an implication for the limits of any observer's view.
A given observer can only directly access facts about systems that have interacted with that observer.
The totality of quantum fact is not, in principle, accessible to any single observer.
Not because some facts are hidden, but because facts are relational, and an observer's relational reach is finite.
There is always a wider relational context that exceeds what any particular observer can register.

That is the precise sense in which there is no view from nowhere.
Not that observers are systematically deceived, or that reality is unknowable, but that every observer occupies a relational standpoint, and from every relational standpoint, there is more to reality than what that standpoint makes available.
The view from everywhere would require an entity in relational contact with everything.
No finite observer is that entity.
Reality is structured so that the whole always exceeds the reach of any particular perspective.

\subsection{A Worked Example: Wigner's Friend}

The abstract statement becomes clearer, and its cost more visible, in a thought experiment the physicist Eugene Wigner devised in the early 1960s.
Rovelli's 1996 paper takes up a version of the same scenario, and it is worth walking through slowly.

Wigner's friend sits inside a sealed laboratory.
She measures the spin of an electron that has been prepared so that, along the vertical axis, it is in an equal superposition of up and down.
Her detector clicks; she reads the display; she writes ``up'' in her notebook.
For her, the matter is settled.
The electron has a definite value, and she knows what it is.

Wigner stands outside the sealed door.
Nothing has passed between him and the laboratory since the experiment began.
If he applies the ordinary equations of quantum mechanics to everything behind the door --- electron, detector, notebook, friend --- those equations never produce a single outcome.
They produce an entangled state of the whole laboratory: one branch in which the electron is up and the notebook reads ``up,'' and one in which the electron is down and the notebook reads ``down,'' held together in superposition.
So is it true, or not, that the friend saw ``up''?

Textbook quantum mechanics has trouble with this question because it says that measurement collapses the state but does not say what counts as a measurement.
If the friend's measurement collapsed the state, Wigner's description is wrong.
If it did not, the friend's certainty seems to be an illusion.
Relational quantum mechanics answers that neither description is wrong, because they are not descriptions of the same relational fact.
Relative to the friend, the electron's spin is up.
Relative to Wigner, there is not yet any value, only a correlation: whatever he eventually finds when he opens the door, the electron and the notebook will agree.
That guaranteed agreement is the consistency constraint at work.

Two features of this example matter for the rest of the chapter.
The first is that Wigner's superposition is not a polite name for his ignorance.
In principle, though far beyond any real laboratory, Wigner could perform an interference measurement on the whole sealed room whose statistics would differ from those of a room that simply contained a definite but unknown result.
The entangled description makes different predictions, which is why the relativity of the outcome cannot be explained away as a gap in Wigner's information.
The second is that nothing in the example depends on the friend being conscious.
Replace her with a photographic plate and the analysis is unchanged.
In Rovelli's usage an ``observer'' is any physical system that interacts with another, and a relative fact is a fact about that interaction.
I will return to this point, because it limits what the convergence can claim.

Ordinary life does not feel this way, and my colleagues and I never disagree about whether the needle on a dial points left or right.
The reason is decoherence~\citep{Zurek2003}.
A macroscopic record, such as the position of a needle, is continuously copied into its environment by scattered light, air molecules, and thermal radiation.
This constant monitoring singles out a small set of robust ``pointer'' states that survive it, a process Zurek calls einselection, and the many copies make the record available to anyone who looks.
Once a fact has been copied that widely, every observer who consults the environment finds the same thing, and the relative fact behaves, for all practical purposes, like an absolute one.
Decoherence does not by itself explain why one outcome rather than another occurs, but it explains why a shared, stable classical world appears within a relational one.
The perspectival structure has not gone away; it has been smoothed over by the sheer number of overlapping perspectives.

I should also be candid that relational quantum mechanics is one interpretation among several.
Hugh Everett's relative-state formulation, QBism, and Bohmian mechanics each read the formalism differently, and Bohmian mechanics in particular restores observer-independent particle positions.
The strong claim of this section, that outcomes are facts only relative to an interacting system, depends on the relational reading.
A weaker claim does not.
On every interpretation, no single measurement on a single system can recover that system's full, unknown quantum state, and no device can copy such a state so as to measure it again~\citep{NielsenChuang2000}.
The quantum world does not hand its whole description to any one interaction.

\section{The Process-Philosophical View: Grades of Experience}
\label{sec:ch9-process}

Whitehead arrived at a similar conclusion from a very different direction.
His starting point was not the quantum measurement problem but a puzzle about experience.
The dominant tradition in Western philosophy since Descartes had divided reality into two radically different kinds of substance: mind and matter.
Mind is characterized by subjectivity, by experience, by the felt quality of being something.
Matter is characterized by extension, by mechanism, by the absence of anything felt.

Whitehead found this division incoherent.
Not because he wanted to eliminate mind and reduce everything to mechanism --- that was the other, materialist horn of the same dilemma.
He found it incoherent because it required the emergence of experience from something entirely devoid of experience, at some unspecifiable threshold.
At what point does matter become aware?
At what level of neural complexity does mechanism become experience?
No one has given a satisfying answer, because the question presupposes a categorical gap that Whitehead thought was the source of the problem rather than a feature of the world.

His solution was to generalize experience downward rather than trying to explain how it emerges upward.
Every actual occasion, at every level of reality, has something like an experiential character --- a perspective, a grasping, a response to its environment.
This does not mean that electrons have rich conscious inner lives.
It means that the most rudimentary physical events exhibit the same formal structure that, at higher levels of integration, becomes what we recognize as experience: taking account of the environment, integrating incoming influences, and settling into a definite character through that process of relational reception.

Whitehead called this process prehension, from the Latin for grasping.
And he argued that different entities differ not in kind but in the depth and richness of their prehensive integration.
A simple physical occasion prehends its environment in a thin, predominantly physical, and largely unconscious way.
A complex living cell prehends its environment with more integration, more internal complexity, more capacity to respond selectively.
A vertebrate nervous system integrates prehensions at an enormously higher level of complexity, introducing something that begins to look like genuine experience.
A conscious human mind integrates at the highest level we know, producing the full richness of conscious thought, feeling, and creativity.

This is Whitehead's hierarchy of grades of experience, and it is structural, not merely taxonomic~\citep{Whitehead1933}.
It is not just a list of things in order of complexity.
It is a claim about the form of reality: that the capacity to take account of the world, to prehend and integrate the environment, is a fundamental feature of every actual entity --- and that this capacity varies, from the minimal in the most elementary physical occasions to the maximal in the most developed forms of consciousness.

Now here is the piece that connects to relational quantum mechanics.
Each actual occasion prehends its environment from its own spatiotemporal standpoint, with its own particular relational resources.
No occasion exhausts what it prehends.
No occasion takes in the full richness of the occasions it grasps.
There is always more to reality than any given prehension can integrate.

In \emph{Adventures of Ideas}, Whitehead is explicit about this~\citep{Whitehead1933}.
The finite perspective is not a defect of knowing; it is a structural feature of being.
To be an actual occasion is to be a finite relational perspective on reality.
The perspective is genuine.
The grasped reality is real.
But the perspective is partial, and its partiality is not an accident or a limitation to be eventually overcome.
It is what it means to be a particular entity rather than everything at once.

Whitehead gives this partiality a precise mechanism in \emph{Process and Reality}~\citep{Whitehead1929}, and two of his technical terms are worth learning.
The first is \emph{objectification}.
When a new occasion prehends a past one, it does not take in that past occasion as it was for itself, in the full immediacy of its own becoming.
It takes it in under a particular aspect, as an object for this new standpoint.
The past occasion's inner life is finished and cannot be relived; what passes forward is a selection of its character, shaped by what the receiving occasion is able to feel.
The second term is \emph{negative prehension}.
To become one definite thing, an occasion must exclude as well as include.
Some of what is available is actively left out of its final synthesis, because it cannot be harmonized with the rest or because the occasion lacks the complexity to hold it.
Perspective, on this account, is achieved partly by elimination.
A standpoint is defined as much by what it declines to feel as by what it feels.

An analogy from human life, and only an analogy: at a string quartet, a trained musician beside me hears voice-leading that I do not register at all.
We hear the same performance, and neither of us is imagining anything, but her capacity for integration lets more of the sound enter her experience as structure, while for me much of it is excluded or felt only as texture.
Whitehead's claim is that something with this form, stripped of consciousness and training, holds for occasions at every level.

One qualification prevents a common misreading.
Whitehead does not hold that rocks, tables, or crystals have experience as unified subjects.
In the development of his thought by later process philosophers, notably David Ray Griffin, a rock is an \emph{aggregational society}: a collection of low-grade occasions with no dominant occasion that unifies the whole into a single perspective~\citep{Griffin2000}.
Only what Griffin, following Charles Hartshorne, calls compound individuals, such as cells and animals, have a unified standpoint above the level of their parts.
This distinction will matter when the \bahai account speaks of what the mineral kingdom ``cannot perceive,'' because process thought would locate the relevant perspective in the mineral's elementary constituents, not in the stone as a whole.

Higher grades of experience do not simply see more.
They integrate more, synthesizing a wider and richer relational field into a more complex and unified act of becoming.
A human mind does not simply have access to more facts than an amoeba.
It integrates its prehensions into a more deeply organized, more internally differentiated, more richly resonant unity.
The quantitative difference in reach becomes a qualitative difference in the character of experience.

Consider what this means for the claim that there is no view from nowhere.
In process ontology, there is no entity that prehends everything without being itself a particular entity with a particular perspective.
God, in Whitehead's framework, has a primordial nature that encompasses the full range of eternal objects --- the abstract forms of possibility --- but even God's consequent nature, the divine engagement with the actual world, takes in the world from God's own perspective and is enriched and shaped by the actual occasions that contribute to it.
The world is present to God in its fullness, but it is present as prehended, as taken in from a relational standpoint.

What Whitehead is insisting on is that the view from nowhere is not just practically unachievable.
It is formally incoherent.
To be an entity capable of knowing is to be an entity with a perspective.
The perspectiveless view would require an entity with no standpoint, no relational location, no particular position in the web of prehensive relations that constitutes reality.
Such an entity would not be a knower.
It would be nothing at all.

The hierarchy of grades is, from this angle, not a hierarchy of approximations to a view from nowhere.
It is a hierarchy of genuine perspectives, each valid at its own level, each disclosing aspects of reality that are inaccessible to the levels below it, and each failing to fully grasp what the levels above it can integrate.
The richest perspective is richer not because it approaches the view from nowhere, but because it integrates more of the world's relational complexity into a genuine, perspectival unity.

This maps onto relational quantum mechanics more cleanly than it might first appear.
Rovelli's different observers correspond, in broad structural terms, to Whitehead's actual occasions at different levels of relational complexity.
Each has a genuine perspective.
Each constitutes genuine facts through its relational interactions.
None exhausts the full reality it engages.
And the notion of a perspective-independent description of quantum reality is, in Rovelli's framework, as incoherent as the notion of a prehension-independent reality in Whitehead's.

\section{The Bahá'í View: The Hierarchy of Kingdoms}
\label{sec:ch9-bahai}

Among the most sustained treatments of observer-dependence in any spiritual tradition, the \bahai writings' account of the hierarchy of kingdoms is striking in its philosophical precision.
\AbdulBaha develops it at length in \emph{Some Answered Questions}, and the argument is not primarily about humility or the limits of human knowing.
It is an ontological argument about the structure of reality~\citep[Part 4]{SAQ}.

The passage deserves to be read carefully.
\AbdulBaha argues that each kingdom of existence --- mineral, vegetable, animal, human, and spiritual --- has its own mode of existence, its own form of participation in reality, its own capacity for perception and relation.
And he argues that each kingdom can, in principle, perceive and understand the kingdoms below it, but cannot perceive or understand the kingdom above it.

The mineral kingdom has its own form of being.
The laws of chemistry and crystallography that govern it are real laws, governing a real mode of existence.
But the mineral has no access to the mode of being that characterizes the vegetable kingdom.
It cannot perceive growth, nutrition, reproduction, or the sensitivity to environment that characterizes plant life.
The mode of being proper to the vegetable kingdom is simply not available within the relational constitution of the mineral.

The vegetable kingdom, in turn, participates in a richer mode of being.
It does what the mineral does, but it also grows, reproduces, and responds to its environment in ways that the mineral cannot.
Yet the vegetable cannot perceive what is proper to the animal kingdom.
Voluntary movement, perception at a distance, the integration of multiple sensory channels into a unified representation of the environment --- these belong to a relational level that the vegetable's constitution does not reach.

The animal kingdom integrates still more.
It does what the vegetable does, and it adds the capacity for locomotion, perception, and something like memory and recognition.
But the animal cannot perceive what is proper to the human kingdom.
Abstract thought, language, mathematics, the capacity to consider possibilities that do not yet exist and to orient behavior toward a future that has not yet arrived --- these belong to a relational complexity that even the most sophisticated non-human animal does not reach.

And the human kingdom, for all its capacities, cannot perceive, in this life, what belongs to the spiritual realm.
The limitations here are structural, not merely quantitative.
The mode of being proper to the spiritual world is simply not available to an entity constituted as human beings in embodied life are constituted.

What is philosophically important here is the directionality of the blindness.
It is not symmetrical.
Higher kingdoms can understand lower kingdoms.
A chemist can fully describe the composition of a mineral.
A human being can understand the vegetable in terms of its chemistry and physiology.
A human being can, to a remarkable degree, understand the animal --- can model its behaviors, its nervous system, its cognition.
But neither the vegetable nor the animal can understand the human.

\AbdulBaha states this asymmetry in so many words in a talk given at Unity Church in Montclair, New Jersey, in May 1912.
He walks up the ladder one rung at a time: the mineral, however far it advances, cannot comprehend the life of the vegetable; the tree, however perfect its growth, cannot share the senses of the animal; the animal, however developed, has no idea of intellect and spirit.
Then he names the principle.
``A lower degree cannot comprehend a higher although all are in the same world of creation,'' and ``Degree is the barrier and limitation''~\citep[Talk 46]{PUP}.
In the same passage he gives the other half of the asymmetry: ``In the human plane of existence we can say we have knowledge of a vegetable, its qualities and product; but the vegetable has no knowledge or comprehension whatever of us''~\citep[Talk 46]{PUP}.

Three features of this passage are easy to miss.
The first is the phrase ``in the same world of creation.''
The kingdoms are not sealed off in different universes.
The rose and the person share one world, stand in the same garden, and interact constantly.
What separates them is not location or causal isolation but degree, which is to say the relational capacity each brings to the encounter.
That is exactly the shape of the claim I drew from the physics: two systems can be in constant interaction and still register quite different facts.

The second feature is where the argument is going.
The talk is not, in the end, about plants and animals.
\AbdulBaha uses the ladder of kingdoms to reach a theological conclusion: if difference of degree already bars comprehension within creation, then the created human being cannot comprehend the essence of the divine Reality, and it is through the Manifestations of God that the divine becomes accessible to human beings at all~\citep[Talk 46]{PUP}.
The hierarchy of kingdoms is thus an argument from the structure of creation to the structure of revelation.
Its top rung is not a higher animal but an unknowable essence, and what reaches the human from that essence reaches it in a form proportioned to human capacity.

The third feature concerns the word ``comprehend.''
When the writings say that the mineral cannot comprehend the vegetable, they are not attributing a thwarted cognitive life to stones.
The word marks a limit of capacity: what a kind of being can take into itself, register, and respond to.
Only on that reading is the comparison with physics and process thought fair.

This asymmetry reflects the structure of the hierarchy itself.
Each level subsumes the capacities of the levels below it while adding capacities irreducible to them, so that the higher perspective is a qualitatively different mode of relational engagement, not a larger pile of the same information.

The \bahai writings also bring a crucial qualification that prevents this hierarchy from collapsing into mere relativism.
The diversity of relational expressions does not mean that reality has no determinate character.
Reality is one.
Its face varies with the relational instrument through which it is perceived.
But the underlying reality is not fragmented or inconsistent.
\Bahaullah, speaking of the powers of the human being, calls them ``one single reality which hath manifold expressions owing to the diversity of its instruments''~\citep[para.~35]{SummonsLordHosts}.
The passage concerns the unity of the human spirit across its many faculties, but the principle it states, one reality and many instruments, is the principle of perspectival realism, not relativism.

The analogy the writings use elsewhere is illuminating: the same sun, many mirrors.
Each mirror reflects the sun truly.
None reflects it completely.
A small, rough mirror reflects a dim and distorted image.
A large, polished mirror reflects a clear and full image.
But both are reflecting the same sun, and no mirror, however polished, exhausts the sun itself.

The differences in the mirrors' reflections are real.
They are not errors.
They are accurate expressions of what the sun is, seen through these particular instruments.
And the fact that no mirror gives a complete reflection does not mean the sun is unknowable.
It means the sun exceeds any particular instrument's capacity.

This is not a claim about human subjectivity.
It is a claim about the architecture of reality.
The hierarchy of kingdoms is a hierarchy of relational instruments, each genuinely disclosing a real aspect of reality, none exhausting it.
The view from everywhere would require an instrument that subsumes all possible modes of being --- an instrument that is not itself at any particular level of the hierarchy, but that somehow includes them all.
That is precisely what the notion of absolute divine knowledge represents in the \bahai writings.
But for any finite mode of existence, reality presents a face proportioned to the instrument.

\section{The Structural Isomorphism: No View From Nowhere}
\label{sec:ch9-convergence}

It would be easy to note that all three accounts say what you observe depends on how you observe, and stop there.
That would miss the deeper structural claim, which is what makes this an isomorphism rather than a family resemblance.
Let me state the common formal point in abstract terms.

There is a set of relational entities.
Each entity has a particular relational constitution --- a particular position in the network of relations that constitutes reality.
From that relational constitution, the entity has access to certain facts about other entities in its environment.
Those facts are genuine facts, not subjective impressions.
They are constituted by the relational interaction between the entity and its environment.
But they are not all the facts.
There is always a wider relational context that exceeds what any particular relational constitution can access.

In relational quantum mechanics, the relational constitution is the observer's physical constitution --- the set of physical systems with which it has interacted and from which it can retrieve information.
In process philosophy, the relational constitution is the actual occasion's prehensive history --- the prior occasions it has integrated and the richness of its integrative capacity.
In the \bahai hierarchy, the relational constitution is the kingdom of existence --- the level of organizational complexity and relational integration that characterizes the entity's mode of being.

The three traditions also agree that none of this is relativism.
All three are perspectival realists: perspectival access is genuine access to real features of the world, and its partiality is a structural feature of being a finite entity in a relational world.

Process thought and the \bahai writings are, in addition, explicitly hierarchical, and hierarchical in the same direction: from narrower to wider relational reach, with the hierarchy extending beyond what present human capacities can grasp.
Whether physics shares that third feature is a harder question, and I take it up below.

The common structure I am identifying is this: a relational ontology generates a plurality of perspectival access, and this plurality is not a defect to be corrected but a constitutive feature of what it means for reality to be relational.
The view from nowhere would be a perspective with no particular relational constitution --- an observer that is not any particular kind of observer.
All three traditions agree that no such entity exists, and that the very concept is incoherent.
What exists is a web of relational perspectives, each genuine, each partial, each constituting its own access to a reality that always exceeds any perspective's reach.

One further structural parallel deserves mention.
In relational quantum mechanics, the incompatibility between two observers' descriptions of the same system is not a contradiction --- it is expected and consistent, as long as the observers have not compared notes.
The incompatibility resolves, not by finding the ``true'' description, but by finding the description that is accurate for the joint system of both observers once they have interacted.
The hierarchical structure of \bahai kingdoms works similarly.
The mineral's description of reality and the human's description of reality are not contradictory.
They are each accurate for the level of relational constitution from which they emerge.
The apparent contradiction resolves not by declaring one wrong, but by recognizing that both are genuine expressions of a reality that exceeds either alone.

This is perspectival realism in its strongest form: the multiplicity of perspectives is not a sign of reality's incoherence but a sign of its inexhaustible depth.

\subsection{The Strongest Objection, and What the Convergence Does Not Claim}

The most serious objection to this convergence comes from the physics itself.
In relational quantum mechanics, perspectives are not ranked.
Any physical system can serve as the reference system for any other.
An electron is as legitimate an ``observer'' of a detector as the detector is of the electron, and a stone is as legitimate an observer as a physicist.
There is no privileged class of observers, no ladder, and no top.
The \bahai hierarchy, by contrast, is ranked and asymmetric: higher kingdoms comprehend lower ones and not the reverse.
So, the objection runs, I have taken a flat, democratic structure from physics and quietly dressed it in the vertical clothing of a religious cosmology.

I think the objection is right about legitimacy and wrong about capacity.
It is true that relational quantum mechanics grants every system an equally valid standpoint.
But it does not grant every system an equal reach.
What one system can register about another is limited by its own physical make-up.
A single two-level system, such as one electron's spin, can carry at most one bit of classical information about anything it interacts with~\citep{NielsenChuang2000}.
It can become correlated with a detector, but it cannot hold a record of the detector's state, compare that record with another, or carry it forward in a way that other systems can consult.
A macroscopic apparatus can do all of these things.
There is therefore a real asymmetry in physics, grounded in physical constitution: richer systems can register more of the relational world than simpler ones, and the simpler ones cannot register the richer ones in anything like their full structure.
That asymmetry is a gradient of capacity, not a ladder of kinds, and I do not want to pretend otherwise.
The hierarchy of kingdoms has its closest counterpart in Whitehead's grades of experience; physics supplies the floor on which both rest, namely that definite facts are relative to a standpoint and that no standpoint takes in the whole.
The isomorphism I claim is in that shared floor, with the gradient of capacity as a genuine but weaker physical parallel to the ranked hierarchy.

It is equally important to say what the convergence does not claim.
It does not claim that consciousness causes the collapse of the wavefunction.
As the Wigner's friend example showed, the relevant observers in relational quantum mechanics need not be minds at all.
It does not claim that quantum mechanics demonstrates the existence of a spiritual kingdom above the human.
Physics is silent on that question, and nothing in this chapter should be read as a physical argument for it.
It does not claim that the \bahai texts contain quantum mechanics in coded form; the texts and the physics share a structure, and the structure is all I am comparing.
And it does not claim that every perspective is as good as every other.
Perspectives are equally real, but they are not equally rich, and the difference between those two statements is the whole difference between perspectival realism and relativism.

\begin{convergencemoment}{Observer-Dependence: What You Can See Depends on What You Are}
In relational quantum mechanics~\citep{Rovelli1996}, a quantum system has no absolute, observer-independent state.
What becomes definite --- what is real in the sense of having a determinate value --- is constituted by the relational interaction between the system and an observer.
Two observers, which need not be minds, may consistently assign different states to the same system if they have not compared notes.
Neither is wrong.
The observer-dependence is not a defect of quantum mechanics; it is a structural feature of relational reality.

In process philosophy~\citep{Whitehead1929,Whitehead1933}, every actual occasion prehends reality from its own relational standpoint, with its own particular prehensive resources.
No occasion exhausts the richness of what it prehends.
Whitehead's hierarchy of grades of experience is a hierarchy of relational perspectives, running from the most rudimentary physical occasions to the most complex forms of conscious integration.
Higher grades integrate more of the world's relational complexity into richer acts of becoming.
None integrates everything.

In the \bahai writings~\citep[Part 4]{SAQ}, the mineral cannot perceive what the vegetable perceives; the vegetable cannot perceive what the animal perceives; the animal cannot perceive what the human perceives; and the human cannot perceive, in this life, what belongs to the spiritual realm.
This hierarchy is not primarily about what different beings know; it is about what different beings are in relation to.
Each level of existence participates in a relational context that exceeds its comprehension.

One structure underlies all three.
What reality discloses depends on what kind of relational entity is doing the disclosing.
There is no view from nowhere --- not because reality is subjective, but because reality is relational, and every perspective on a relational world is itself a relational position.
The partiality of perspective is not a defect.
It is what it means to be a finite entity in an inexhaustible world.
\end{convergencemoment}
